\clearpage
\VerdierPhysicalPageAnchor{0108}%
\begingroup
\small
\setlength{\abovedisplayskip}{4pt plus 2pt minus 2pt}%
\setlength{\belowdisplayskip}{4pt plus 2pt minus 2pt}%

\VerdierSemanticAnchor{ch1-app-signs}
\section*{付録：符号の選択に関する注記}
\addcontentsline{toc}{section}{付録：符号の選択に関する注記}

\VerdierSemanticAnchor{ch1-app-signs-cone}
写像錐関手は、その随伴性 (3.1.3) により一意な同型を除いて定まる。
$\operatorname{tr}_1(C)$ の射に $\operatorname{comp}(C)$ の三角を容易に対応させる
ため、シフト関手として「射 $X\longrightarrow 0$ の写像錐」を選んだ
（(3.1.7) および (3.2.2.7)）。すなわち、次数をずらし、微分の符号を変える。

\VerdierSemanticAnchor{ch1-app-signs-strict}
$F\colon\prod_{i\in[n]}C_i\longrightarrow C$ を多加法的関手とする。圏
$\prod_{i\in[n]}\operatorname{comp}(C_i)$ は自然に $\mathbb Z[n]$-圏の構造をもつ。
関手
\[
  F^{[n]}\colon\prod_{i\in[n]}\operatorname{comp}(C_i)
  \longrightarrow\operatorname{comp}^{n}(C)
\]
を厳密な $\mathbb Z[n]$-関手にしようとすると、われわれが採用したシフトの
規約、すなわち一つの方向に次数をずらし、対応する微分の符号を変えるという
選択に導かれる。

\VerdierSemanticAnchor{ch1-app-signs-simple-complex}
付随する単純複体関手の選択には [1] の規約を採用した（ただし、われわれの
多重複体では微分どうしは可換である）。これにより関手
$K_{\Pi}F\colon\prod_i K(C_i)\longrightarrow K(C)$ を定義できる。このとき
$K_{\Pi}F(X_1,\ldots,X_i[1],\ldots,X_n)$ は
$K_{\Pi}F(X_1,\ldots,X_i,\ldots,X_n)[1]$ と異なる。それでも関手
$K_{\Pi}F$ が完全三角を完全三角に送ると述べるために、同型

\VerdierSemanticAnchor{ch1-app-signs-shift-isomorphism}
\[
  p_i\colon K_{\Pi}F(X_1,\ldots,X_i[1],\ldots,X_n)
  \xrightarrow{\ \sim\ }
  K_{\Pi}F(X_1,\ldots,X_i,\ldots,X_n)[1]
\]
を定義することになる（(3.4.3) の証明）。

\VerdierSemanticAnchor{ch1-app-signs-anticommutativity}
すると $i\ne j$ のとき、次の反可換図式が得られる。

\VerdierSemanticAnchor{ch1-app-signs-diagram}
\begin{center}
  \resizebox{\textwidth}{!}{%
    $\begin{tikzcd}[ampersand replacement=\&,column sep=3.2em,row sep=4.3em]
      K_{\Pi}F(X_1,\ldots,X_i[1],\ldots,X_j[1],\ldots)
        \arrow[r,"p_j"]
        \arrow[d,"p_i"']
        \& K_{\Pi}F(X_1,\ldots,X_i[1],\ldots,X_j,\ldots)[1]
          \arrow[d,"p_i"] \\
      K_{\Pi}F(X_1,\ldots,X_i,\ldots,X_j[1],\ldots)[1]
        \arrow[r,"p_j"]
        \& K_{\Pi}F(X_1,\ldots,X_i,\ldots,X_j,\ldots)[2].
    \end{tikzcd}$%
  }
\end{center}

\VerdierSemanticAnchor{ch1-app-signs-koszul-choice}
続いて (1.5.2) では、これらの同型を同型 $p(\sigma)$ へ整合的に拡張する
方法を示した。この同型は各 $\sigma\in\mathbb Z[n]$ に対して定義される。
ここには一つの
任意な選択が入る。それは、それ自身はよく定まったコホモロジー類、すなわち
Koszul 類 (1.6.3) の中で余サイクルを選ぶことである。一般に受け入れられている
規則（例えば [1]）に合わせるため、すべての場合に Koszul 余サイクル (1.6.6)
を選んだ。
\endgroup
